\section{Acoplamento e Pseudo Algoritmo}

Os domínios são acoplados por meio de interpolações e condições de contorno, da seguinte forma:

O escoamento principal recebe informações do filme através de condições de contorno nas paredes que são compartilhadas por ambos os domínios, as quais são detalhadas na \ref{sec:condicoes_de_contorno}. 

Enquanto isso, o filme recebe informações do escoamento principal de duas formas: os campos de velocidade vem como condição de contorno nas paredes internas do filme, sendo paredes internas aquelas que não coincidem nem com uma fronteira imersa nem com o domínio físico do escoamento, através dos níveis intermediários e o campo de pressão é interpolado diretamente do \textit{coreflow} para o filme, sendo que essa propriedade não é resolvida neste domínio.

Para que os níveis intermediários sejam capazes de fornecer as condições de contorno, os campos de velocidade são interpolados do escoamento principal utilizando interpoladores conservativos.

convencionando o tempo precedente como $t^n$ e o tempo seguinte como $t^{n+1}$, temos que a sequência de passos para o método em um determinado passo de tempo começa com as condições do filme e do gás em $t^n$, seguindo a seguinte sequência de etapas:

\begin{enumerate}
    \item Calcula-se o passo de tempo do \textit{core flow}, $\Delta t_\mathrm{core}$.
    \item Calcula-se o passo de tempo do filme, $\Delta t_\mathrm{filme}$.
    \item Calcula-se o número de subciclos, arredondando para cima a razão entre os passos de tempo: $N_\mathrm{passos} = \left\lceil \Delta t_\mathrm{core} / \Delta t_\mathrm{filme} \right\rceil$.
    \item Define-se o passo de tempo efetivo do filme, $\delta t_\mathrm{filme} = \Delta t_\mathrm{core} / N_\mathrm{passos}$, de modo que $\delta t_\mathrm{filme} \le \Delta t_\mathrm{filme}$ e $N_\mathrm{passos}\,\delta t_\mathrm{filme} = \Delta t_\mathrm{core}$.
    \item O \textit{core flow} avança um passo de tempo $\Delta t_\mathrm{core}$ usando as condições de contorno do filme em $t^{n}$, obtendo os campos de velocidade e pressão em $t^{n+1}$.
    \item O campo de velocidade do \textit{core flow} é interpolado para os níveis intermediários $t^{n,k} = t^{n} + k\,\delta t_\mathrm{filme}$, com $k = 1, \dots, N_\mathrm{passos}$.
    \item O campo de pressão do \textit{core flow} é interpolado até o ltop passando pelos níveis intermediários $t^{n,k}$.
    \item O filme avança $N_\mathrm{passos}$ subciclos de tamanho $\delta t_\mathrm{filme}$, usando em cada nível $t^{n,k}$ os campos interpolados do \textit{core flow} como condição de contorno, até obter o campo de velocidade do filme em $t^{n+1}$.
    \item O processo se repete para o próximo passo de tempo do \textit{core flow}.
\end{enumerate}

Com o seguinte pseudo algoritmo equivalente:

\begin{algorithm}[H]
\caption{Sequência do algoritmo com subciclagem temporal do filme.}
\label{alg:sequencia-subciclagem}
\begin{algorithmic}[1]
\Require $u^0$, $p^0$, $u_\mathrm{filme}^0$, $t^0$, $t_\mathrm{final}$
\State $n \gets 0$
\While{$t^n < t_\mathrm{final}$}
  \State Calcular $\Delta t_\mathrm{core}$
  \State Calcular $\Delta t_\mathrm{filme}$
  \State $N_\mathrm{passos} \gets \lceil \Delta t_\mathrm{core} / \Delta t_\mathrm{filme} \rceil$
  \State $\delta t_\mathrm{filme} \gets \Delta t_\mathrm{core} / N_\mathrm{passos}$
  \State $(u^{n+1}, p^{n+1}) \gets \textsc{AvançaCore}\big(u^{n}, p^{n}, \mathcal{B}_\mathrm{filme}^{n}, \Delta t_\mathrm{core}\big)$
  \For{$k = 1, \dots, N_\mathrm{passos}$}
    \State $t^{n,k} \gets t^n + k\,\delta t_\mathrm{filme}$
    \State $u^{n,k} \gets \textsc{Interpola}\big(u^{n}, u^{n+1}, t^{n,k}\big)$
    \State $p^{n,k} \gets \textsc{Interpola}\big(p^{n}, p^{n+1}, t^{n,k}\big)$
  \EndFor
  \State $u_\mathrm{filme}^{n,0} \gets u_\mathrm{filme}^{n}$
  \For{$k = 1, \dots, N_\mathrm{passos}$}
    \State $u_\mathrm{filme}^{n,k} \gets \textsc{AvançaFilme}\big(u_\mathrm{filme}^{n,k-1}, u^{n,k}, p^{n,k}, \delta t_\mathrm{filme}\big)$
  \EndFor
  \State $u_\mathrm{filme}^{n+1} \gets u_\mathrm{filme}^{n,N_\mathrm{passos}}$
  \State $t^{n+1} \gets t^n + \Delta t_\mathrm{core}$, \quad $n \gets n+1$
\EndWhile
\end{algorithmic}
\end{algorithm}

Sendo, nessas representações, $u$ representa o campo de velocidades, $t$ o tempo, $p$ a pressão e $N_{passos}$ o número de passos 
